\documentclass[11pt]{article}
\usepackage[margin=1in]{geometry}
\usepackage{booktabs}
\usepackage{array}
\usepackage{amsmath}
\usepackage[table]{xcolor}

\title{ENSC 453/894 --- Lab 5: CUDA Image Blur Optimization Report}
\author{}
\date{}

\begin{document}
\maketitle

\noindent\textbf{GPU:} Nvidia Ampere A4000 (48 SMs, 128 cores/SM, 164\,KB shared mem/SM, 2048 threads/SM)\\
\textbf{Image:} 3840 $\times$ 2160 (grayscale, single-channel float)\\
\textbf{BLUR\_SIZE:} 21 (blur window = 43 $\times$ 43 = 1,849 neighbours per pixel)\\
\textbf{Kernel repeats:} 10 per run. All times include data transfer.

%======================================================================
\section{Shared Memory Optimization}
%======================================================================

\subsection{Strategy --- Separable Row/Column Sum Reuse}

The baseline kernel performs $O((2R{+}1)^2) = O(1{,}849)$ global memory reads per output
pixel. The optimized kernel exploits the fact that vertically adjacent pixels share
42 out of 43 rows of their blur window. It proceeds in three phases on shared memory:

\begin{enumerate}[]
  \item \textbf{Tile Load:} Each block cooperatively loads a tile of
        $(\text{OUTPUT\_DIM} + 2 \times \text{BLUR\_SIZE})^2$ pixels from global memory
        into shared memory, zero-padding out-of-bounds pixels.
  \item \textbf{Row Sums:} For every tile row and every output-column position,
        accumulate 43 consecutive elements horizontally. Each row sum is reused by
        up to 43 output pixels that share the same row in their blur window.
  \item \textbf{Column Accumulation:} For each output pixel, sum 43 row sums
        vertically and divide by the analytically computed neighbour count:
        $(x_{\text{end}} - x_{\text{start}} + 1) \times (y_{\text{end}} - y_{\text{start}} + 1)$.
\end{enumerate}

This reduces per-pixel work from $O(1{,}849)$ to $O(2 \times 43) = O(86)$---a
\textbf{21.5$\times$ arithmetic reduction}.

\subsection{Tile Size Analysis}

The tile multiplier $I$ (\texttt{TILE\_I}) controls how many output pixels each block
produces. Each block of $16 \times 16 = 256$ threads outputs an $(I \cdot 16)^2$
region; each thread handles $I^2$ pixels.

\begin{table}[h]
\centering
\caption{Tile size sweep (shared memory kernel, 16$\times$16 blocks)}
\smallskip
\begin{tabular}{c c c r c c c}
\toprule
\textbf{TILE\_I} & \textbf{Output} & \textbf{Tile} & \textbf{Shared Mem} & \textbf{Blocks/SM} & \textbf{Occupancy} & \textbf{Time (s)} \\
\midrule
1 & 16$\times$16 & 58$\times$58   & 16.8\,KB & 8 (thread-lim.) & 100\%  & 0.0373 \\
2 & 32$\times$32 & 74$\times$74   & 30.6\,KB & 5 (smem-lim.)   & 62.5\% & 0.0317 \\
\rowcolor{green!15}
\textbf{3} & \textbf{48$\times$48} & \textbf{90$\times$90} & \textbf{48.5\,KB} & \textbf{3 (smem-lim.)} & \textbf{37.5\%} & \textbf{0.0300} \\
4 & 64$\times$64 & 106$\times$106 & 70.4\,KB & 2 (smem-lim.)   & 25\%   & 0.0308 \\
\bottomrule
\end{tabular}
\end{table}

\textbf{Best: TILE\_I\,=\,3} (0.030\,s). Despite only 37.5\% occupancy, the larger
output region amortizes halo overhead: the ratio of useful output pixels to total tile
pixels is $48^2/90^2 = 28.4\%$, versus $16^2/58^2 = 7.6\%$ for $I{=}1$. At $I{=}4$
the halo ratio improves further (36.5\%) but occupancy drops to 25\%, and the SM can
no longer hide memory latency with so few active warps, causing performance to plateau.

%======================================================================
\section{Thread Block Size Analysis}
%======================================================================

Tested on the naive kernel (global memory, no tiling) with pinned-memory transfers
and \texttt{-O3 -use\_fast\_math}:

\begin{table}[h]
\centering
\caption{Block size sweep (naive kernel, pinned memory)}
\smallskip
\begin{tabular}{c c c c c}
\toprule
\textbf{Block Size} & \textbf{Threads/Block} & \textbf{Blocks/SM} & \textbf{Occupancy} & \textbf{Time (s)} \\
\midrule
8$\times$8   & 64   & 32 (max-block lim.) & 100\% & 0.2184 \\
\rowcolor{green!15}
\textbf{16$\times$16} & \textbf{256} & \textbf{8} & \textbf{100\%} & \textbf{0.2152} \\
32$\times$32 & 1024 & 2                   & 100\% & 0.2185 \\
32$\times$8  & 256  & 8                   & 100\% & 0.2155 \\
16$\times$8  & 128  & 16                  & 100\% & 0.2160 \\
\bottomrule
\end{tabular}
\end{table}

All configurations achieve 100\% occupancy. \textbf{16$\times$16 is fastest} because:
\begin{itemize}[]
  \item The square shape maximises 2D spatial locality for the blur window.
  \item 256 threads (8 warps) provides good latency hiding with minimal register
        pressure, unlike the 1024-thread 32$\times$32 configuration.
  \item 8$\times$8 blocks hit the 32-block/SM hardware limit, adding scheduling overhead.
\end{itemize}

%======================================================================
\section{CPU--GPU Data Transfer Optimization}
%======================================================================

\subsection{Changes from Baseline}
\begin{enumerate}[]
  \item \texttt{cudaFree(0)} warmup --- forces CUDA context initialization before
        timing begins (the baseline includes ${\sim}$1\,s of context init +
        \texttt{cudaMalloc} in its timed section).
  \item \textbf{Pinned host memory} --- \texttt{cudaMallocHost} replaces pageable
        buffers, enabling direct DMA transfers over PCIe.
  \item \textbf{Compiler flags} --- \texttt{-O3 -use\_fast\_math} added to \texttt{nvcc}
        (baseline had no optimization flags).
\end{enumerate}

\subsection{Results}

\begin{table}[h]
\centering
\caption{Data transfer optimization (naive kernel, 16$\times$16 blocks)}
\smallskip
\begin{tabular}{l c c}
\toprule
\textbf{Configuration} & \textbf{Time (s)} & \textbf{Speedup} \\
\midrule
Baseline (pageable mem, no warmup, no \texttt{-O3}) & 1.246 & 1.00$\times$ \\
\rowcolor{green!15}
Data Transfer Optimized (pinned, warmup, \texttt{-O3}) & 0.215 & 5.79$\times$ \\
\bottomrule
\end{tabular}
\end{table}

The 5.79$\times$ speedup comes primarily from excluding CUDA context initialization
from timing (${\sim}$1\,s) and from faster DMA transfers with pinned memory. The
kernel itself is unchanged.

%======================================================================
\section{Control Divergence Analysis}
%======================================================================

\subsection{Naive Kernel (16$\times$16 blocks)}

\noindent\textbf{Grid:} $135 \times 240 = 32{,}400$ blocks, each with 8 warps
$\Rightarrow$ \textbf{259,200 total warps.}

\medskip\noindent\textbf{Source 1 --- Outer boundary check}
(\texttt{if (x >= width || y >= height) return}):
The grid perfectly tiles the image ($135 \times 16 = 2{,}160$;
$240 \times 16 = 3{,}840$). \textbf{Zero divergent warps.}

\medskip\noindent\textbf{Source 2 --- Inner-loop clamping}
(\texttt{if (nx >= 0 \&\& nx < width \&\& ny >= 0 \&\& ny < height)}):
Threads within 21 pixels of any image edge take a different branch from interior
threads on some loop iterations.

\begin{itemize}[]
  \item Divergent block columns: 0--1 (left, $x < 21$) and 133--134 (right, $x > 2138$) = 4
  \item Divergent block rows: 0--1 (top) and 238--239 (bottom) = 4
  \item Interior blocks: $(135 - 4) \times (240 - 4) = 30{,}916$
  \item Divergent blocks: $32{,}400 - 30{,}916 = 1{,}484$
  \item Upper-bound divergent warps: $1{,}484 \times 8 = 11{,}872$
\end{itemize}

\[
  \text{Divergence} = \frac{11{,}872}{259{,}200} = \mathbf{4.58\%}
\]

\subsection{Shared Memory Kernel (TILE\_I = 2, 16$\times$16 blocks)}

\noindent\textbf{Grid:} $68 \times 120 = 8{,}160$ blocks $\Rightarrow$
\textbf{65,280 total warps.}

\begin{table}[h]
\centering
\caption{Divergence by kernel phase (shared memory kernel)}
\smallskip
\begin{tabular}{l l l}
\toprule
\textbf{Phase} & \textbf{Conditionals} & \textbf{Divergence} \\
\midrule
Phase 1 --- Tile Load & Ternary at image edges & ${\sim}$4.6\% of blocks (predicated) \\
Phase 2 --- Row Sums  & Fixed loop $k = 0\ldots42$, none & 0\% \\
Phase 3 --- Col Sums  & Boundary check, warp-uniform & ${\sim}$0\% \\
\bottomrule
\end{tabular}
\end{table}

\textbf{Effective computation divergence $\approx$ 0\%.} The separable design
eliminates all data-dependent branches from the inner loops. Only tile loading at
image edges incurs minor predicated divergence with negligible cost.

%======================================================================
\section{Overall Performance Summary}
%======================================================================

\begin{table}[h]
\centering
\caption{End-to-end comparison across branches}
\smallskip
\begin{tabular}{l l l c c}
\toprule
\textbf{Branch} & \textbf{Kernel} & \textbf{Transfer} & \textbf{Time (s)} & \textbf{Speedup} \\
\midrule
\texttt{baseline}             & Naive (global mem)     & Pageable       & 1.246 & 1.00$\times$ \\
\texttt{data\_transfer-OPUS}  & Naive (global mem)     & Pinned+warmup  & 0.215 & 5.79$\times$ \\
\rowcolor{green!15}
\texttt{shared\_memory-OPUS}  & Shared mem + sep.\ reuse & Pinned+warmup & 0.030 & \textbf{41.5$\times$} \\
\bottomrule
\end{tabular}
\end{table}

\subsection*{Speedup Breakdown}

\begin{center}
\begin{tabular}{ccccc}
Baseline & $\xrightarrow{\;5.79\times\;}$ & Data Transfer & $\xrightarrow{\;7.17\times\;}$ & Shared Memory \\
1.246\,s & & 0.215\,s & & 0.030\,s \\
& {\footnotesize pinned mem + warmup} & & {\footnotesize separable reuse} &
\end{tabular}
\end{center}

\noindent The shared memory kernel's 7.17$\times$ speedup over the naive kernel
(both with pinned memory) comes from:

\begin{enumerate}[]
  \item \textbf{Arithmetic reduction:} $O(86)$ vs $O(1{,}849)$ operations per pixel---a
        21.5$\times$ reduction in computation.
  \item \textbf{Memory hierarchy:} All intermediate reads from shared memory
        (${\sim}$19\,TB/s) vs global memory (${\sim}$448\,GB/s)---a 42$\times$
        bandwidth advantage.
\end{enumerate}

\noindent\textbf{Best configuration:} TILE\_I\,=\,3, BLOCK\_SIZE\,=\,16$\times$16,
achieving \textbf{0.030\,s} total on the Nvidia A4000.

\end{document}
